\chapter{有限体}
    \section{素数位数の有限体の構成法}
        $p$を素数とし，$S\coloneq\{0,1,\dotsc,p-1\}$とする。$S$の任意の元$a,b$に対する演算$+'$と$*'$を次式で定義する。
        \begin{itemize}
            \item $a+'b = (a+b)\% p$
            \item $a*'b = (ab)\% p$
        \end{itemize}
        このようにすると代数構造$(S,+',*')$が体を成すことを説明する。
        \par
        まず$+'$に関する結合法則, 単位元と唯一の逆元の存在，$*'$に関する結合法則と単位元の存在，また，$*'$が$+'$の上に分配的であることは合同式の基本的な知識があれば明らかであろう。
        あとは，$*'$に関する唯一の逆元が存在すること，すなわち任意の$0 \neq a\in S$に対して唯一の$b\in S$が存在して$a*'b=1$となることを示せばよいのだが，\ref{a,mが互いに素のとき，0,a,...a(m-1)をmで割った余りは全て異なる}より，それが保証される。
        \par
        有限体はGalois体とも呼ばれ，位数$p$の有限体を$\GaloisField{p}$と表す。
    \section{位数が素数の冪乗である有限体の構成法}
        \subsection{既約多項式を用いる方法}
            $p$を素数とする。$\GaloisField{p}$上の既約多項式を任意に一つ選んで$P$とし，$m = \deg(P)$とする。
            $S$を，$\GaloisField{p}$上の次数$m-1$以下の全ての多項式の集合を$S$とする($|S|=p^m$)。
            $S$の任意の元$a,b$に対する演算$+'$と$*'$を次式で定義する。
            \begin{itemize}
                \item $a+'b = (a+b)\% P$
                \item $a*'b = (ab)\% P$
            \end{itemize}
            このようにすると代数構造$(S,+',*')$が体を成すことを説明する。
            \par
            まず$+'$に関する結合法則, 単位元と唯一の逆元の存在，$*'$に関する結合法則と単位元の存在，また，$*'$が$+'$の上に分配的であることは合同式, 素数位数の有限体の基本的な知識があれば明らかであろう。
            あとは，$*'$に関する唯一の逆元が存在すること，すなわち任意の$0 \neq a\in S$に対して唯一の$b\in S$が存在して$a*'b=1$となることを示せばよい。
            少なくとも一つ存在することはB\'ezoutの等式により保証される。
            唯一であることを背理法で示す。
            仮に異なる$b_1,b_2 \in S$が存在して$a*'b_1 = a*'b_2 = 1$であるとすると$a*'(b_1 - b_2) = 0$であるが，これは$P$は既約であることと矛盾する。
            なぜならば$a,b_1,b_2$は$P$と互いに素であるから，$P \nmid a*'(b_1 - b_2)$であるからである。
        \subsection{既約多項式の根を用いる方法(根体)}
            $p$を素数とする。
            $\GaloisField{p}$上の既約多項式を任意に一つ選んで$P$とし，$m = \deg(P)$とする。
            $\alpha$を$P$の任意の根とする。
            $S \coloneq \setComprehension*{\sum_{i=0}^{m-1}a_i\alpha^i}{a_0,\dotsc,a_{m-1} \in \GaloisField{p}}$とする。
            $S$上の演算として$\GaloisField{p}$上の多項式に対する加法と乗法をそのまま使うと，$S$とこの演算の組が体を成すことを説明する。
            \par
            まず加法に関する結合法則, 単位元と唯一の逆元の存在，乗法に関する結合法則と単位元の存在，また，乗法が加法の上に分配的であることは$P(\alpha) = 0$であること, 素数位数の有限体, 多項式の除法の基本的な知識があれば明らかであろう。
            あとは，乗法に関する唯一の逆元が存在すること，すなわち任意の$0 \neq a\in S$に対して唯一の$b\in S$が存在して$ab=1$となることを示せばよい。
            少なくとも一つ存在することは，$a$に現れる$\alpha$を不定元$x$で置換して多項式に戻してからB\'ezoutの等式を考えることで理解できる。
            唯一であることは背理法で示せる。
            仮に異なる$b_1,b_2 \in S$が存在して$ab_1 = ab_2 = 1$であるとすると$a(b_1 - b_2) = 0$であるが，一方で$a, b_1 - b_2 \neq 0$だから$a(b_1 - b_2) \neq 0$であり，矛盾している。
    \section{原始多項式}
        \subsection{定義}
            \cite{prim_poly_Wikipedia} に述べられているが，$\GaloisField{p}$ 上の $m\in\naturalNumbers$ 次の既約多項式 $p(x)$ が原始多項式であることは，$x^n-1\;(n\in\naturalNumbers)$ が $p(x)$ で割り切れる最小の $n$ が $p^m-1$ であることと同値である。
        \subsection{原始多項式と生成元}
            \begin{shadebox}
                次の命題は真である。\\
                $\GaloisField{p}$ 上の $m\in\naturalNumbers$ 次の既約多項式 $p(x)$ が原始多項式である $\iff\quad p(x)$ を既約多項式に選んで構成された $\GaloisField{p^m}$ に於いて $x$ が生成元である
            \end{shadebox}
            \begin{proof}
                \quad\par\noindent
                ($\Rightarrow$)
                \par
                集合 $S_0\coloneq\{x,x^2,x^3,\dots,x^{p^m-2}\}$ の元が 0, 1 いずれとも合同でなく，かつ互いに異なることを示す。
                $S_0$ の元が 1 でないことは原始多項式の定義 (<<原始多項式>>) から明らかである。
                \par
                $S_0$ の元が 0 でないことは背理法で示せる。
                仮にある $k\in\{1,2,\dots,p^m-2\}$ について $x^k\equiv 0$ なら $x^l\equiv 0\;(l\in\naturalNumbers,l\geq k)$ であるから $x^{p^m-1}\equiv 1$ に矛盾する。
                \par
                $S_0$ の元が互いに異なることも背理法で示せる。
                仮に $a,b\in\naturalNumbers,p^m-2\geq b>a\geq1$ が存在して $x^b\equiv x^a$ であるとすると $x^a(x^{b-a}-1)\equiv 0$ である。
                原始多項式の定義から $x^{b-a}-1\not\equiv 0$ なので $x^a\equiv 0$ であるが，これは先程示した「$S_0$ の元が 0 でない」に矛盾する。
                \par
                以上より $\{0,1\}\cup S_0$ の要素は互いに異なり，要素数は $p^m-1$ であるから $x$ は $\GaloisField{p^m}$ の生成元である。
                \par\noindent
                ($\Leftarrow$)
                \par
                前提から集合 $S\coloneq\{0,1,x,x^2,x^3,\dots,x^{p^m-2}\}$ の要素は互いに異なる（☆1）。
                \par
                まず $x^{p^m-1}\not\equiv 0$ を背理法で示す。
                仮に $x^{p^m-1}\equiv 0$ とすると $x^{p^m-2}x\equiv 0$ であり、☆1 から $x^{p^m-2}\not\equiv 0$ なので $x\equiv 0$ となり☆1 に矛盾する。
                \par
                次に $x^{p^m-1}\not\in S\setminus\{0,1\}$ を背理法で示す。
                仮にある $a\in\{1,2,\dots,p^m-2\}$ が存在して $x^{p^m-1}\equiv x^a$ であるとすると、 $x^a(x^{p^m-1-a}-1)\equiv 0$ であるが☆1 から $x^{p^m-1-a}-1\not\equiv 0$ なので $x^a\equiv 0$ となり☆1 に矛盾する。
                \par
                $x$ が $\GaloisField{p^m}$ の生成元であるから $x^{p^m-1}$ は必ず $S$ の要素のいずれかと合同である。
                このことと、以上に示したことから $x^{p^m-1}\equiv 1$ である。
            \end{proof}
        \subsection{\texorpdfstring{$\sum_{i=0}^n a_ix^i$}{a0 x0 + a1 x1 + ...}が原始多項式ならば\texorpdfstring{$\sum_{i=0}^n a_{n-i}x^i$}{a\{n-1\} x0 + a\{n-2\} x1 + ...}も原始多項式である}
            \begin{shadebox}
                $p$を素数とする。
                $n$次多項式$f_1(x) \coloneq \sum_{i=0}^n a_ix^i \in \GaloisField{p}[x]$が原始多項式であれば，係数の順番を逆にした多項式$f_2(x) \coloneq \sum_{i=0}^n a_{n-i}x^i$もまた原始多項式である。
            \end{shadebox}
            \begin{proof}
                \quad\par
                \ref{既約多項式の係数の順番反転}より$f_2(x)$は$\GaloisField{p}$上で既約である。
                あとは$f_2(x) \mid x^l - 1$となる最小の自然数$l$が$p^m-1$であることを示せば良い。
                \par
                まず$f_2(x) \mid x^{p^n-1}-1$を示す。
                $f_1(x)$が原始多項式だから適当な$p^n-1-n$次多項式$g_1(x) \in \GaloisField{p}[[x]]$が存在して次式が成り立つ。
                \[ x^{p^n-1} - 1 = g_1(x)f_1(x) \]
                上式をLaurent多項式に拡張して両辺の$x$を$x^{-1}$で置き換えると
                \[ \frac{1}{x^{p^n-1}} - 1 = g_1(x^{-1})f_1(x^{-1}) \]
                両辺に$x^{p^n-1}$を掛けると
                \[ 1 - x^{p^n-1} = x^{p^n-1-n}g_1(x^{-1}) x^n f_1(x^{-1}) =: g_2(x)f_2(x) \]
                \[ \therefore x^{p^n-1} -1 \equiv 0 \mod f_2(x) \]
                ここに$g_2(x)$は$g_1(x)$の係数の順番を逆にしたものである。
                \par
                次に$l=1,2,\dotsc,p^n-2$について$f_2(x) \nmid x^l-1$を背理法で示す。
                仮にある$l$に対して$l-n$次多項式$h_1(x) \in \GaloisField{p}[x]$が存在して
                \[ x^l - 1 = h_1(x)f_2(x) \]
                が成り立つとする。
                上式をLaurent多項式に拡張して両辺の$x$を$x^{-1}$で置き換えると
                \[ \frac{1}{x^l} - 1 = h_1(x^{-1})f_2(x^{-1}) \]
                両辺に$x^l$を掛けると
                \[ 1 - x^l = x^{l-n}h_1(x^{-1}) x^n f_2(x^{-1}) =: h_2(x)f_1(x) \]
                \[ \therefore x^l - 1 \equiv 0 \mod f_1(x) \]
                ここに$h_2(x)$は$h_1(x)$の係数の順番を逆にしたものである。
                これは$f_1(x)$が原始多項式であることに矛盾する。
                \par
                以上より$f_2(x) \mid x^l - 1$となる最小の自然数$l$は$p^m-1$である。
            \end{proof}